% Part: many-valued-logic
% Chapter: syntax-and-semantics
% Section: introduction

\documentclass[../../../include/open-logic-section]{subfiles}

\begin{document}

\olfileid[pt]{mvl}{syn}{int}

\olsection{Introdução}

Na lógica clássica, lidamos com !!{formula}s construídas a partir de
!!{propositional variable}s por meio dos conectivos proposicionais
$\lnot$, $\land$, $\lor$, $\lif$ e $\liff$. Quando definimos uma
semântica para a lógica clássica, fazemos isso usando os dois valores de verdade
$\True$ e $\False$. Interpretamos as !!{propositional variable}s em
!!a{valuation}~$\pAssign{v}$, que atribui esses valores de verdade
$\True$, $\False$ às !!{propositional variable}s. Toda
!!{valuation} determina então um valor de verdade $\pValue{v}(!A)$ para qualquer
!!{formula}~$!A$, e !!^a{formula} é satisfeita em
!!a{valuation}~$\pAssign{v}$, $\pSat{v}{!A}$, se, e somente se, $\pValue{v}(!A) =
\True$.

As lógicas multivaloradas são generalizações da lógica clássica bivalorada
que admitem mais valores de verdade que apenas $\True$ e $\False$. Assim, na
lógica multivalorada, !!a{valuation}~$\pAssign{v}$ é uma função que
atribui a cada !!{propositional variable}~$p$ um dentre uma coleção de
possíveis valores de verdade. Em geral, chamaremos de~$V$ o conjunto dos valores de verdade
permitidos. A lógica clássica é uma lógica multivalorada na qual $V =
\{\True, \False\}$, e o valor de verdade~$\pValue{v}(!A)$ é calculado
usando as conhecidas tabelas-verdade características dos conectivos.

Quando acrescentamos valores de verdade, há mais de uma opção natural
para calcular~$\pValue{v}(!A)$ para os conectivos que lemos
como ``e'', ``ou'', ``não'' e ``se---então''. Portanto, uma lógica multivalorada
é determinada não apenas pelo conjunto de valores de verdade, mas também pelas
\emph{funções de verdade} que decidimos usar para cada conectivo. Uma vez
selecionadas essas funções para uma lógica multivalorada~$\Log L$, contudo, o
valor de verdade $\pValue{v}(!A)[\Log L]$ é unicamente determinado pela
valoração, assim como na lógica clássica. As lógicas multivaloradas, a exemplo da
lógica clássica, são \emph{funcionais em relação à verdade}.

Com esses elementos semânticos em mãos, podemos passar a definir os
análogos dos conceitos semânticos de tautologia, consequência e
satisfatibilidade. Na lógica clássica, !!a{formula} é uma tautologia se
seu valor de verdade $\pValue{v}(!A) = \True$ para qualquer~$\pAssign{v}$. Na
lógica multivalorada, também precisamos generalizar um pouco essa ideia. Em primeiro
lugar, não se exige que o conjunto de valores de verdade~$V$
contenha~$\True$. Por exemplo, algumas lógicas multivaloradas usam números,
como todos os números racionais entre $0$ e~$1$, como conjunto de valores de verdade.
Nesse caso, $1$~costuma desempenhar o papel de~$\True$. Em
outras lógicas, não apenas um, mas vários valores de verdade o fazem. Por isso,
exigimos que toda lógica multivalorada tenha um conjunto~$V^+$ de
\emph{valores designados}. Podemos então dizer que !!a{formula} é
satisfeita em !!a{valuation}~$\pAssign{v}$, $\pSat{v}{!A}[\Log L]$, se, e somente se,
$\pValue{v}(!A)[\Log L] \in V^+$. !!^a{formula}~$!A$ é uma tautologia da
lógica, $\Entails[\Log L] !A$, se, e somente se, $\pValue{v}(!A) \in V^+$ para
qualquer~$\pAssign{v}$. E, por fim, dizemos que $!A$ é consequência de um conjunto
de !!{formula}s, $\Gamma \Entails[\Log L] !A$, se toda !!{valuation} que
satisfaz todas as !!{formula}s de~$\Gamma$ também satisfaz~$!A$.

\end{document}
